\hypertarget{appendix-c-a-type-theory-primer}{%
\chapter{A Type Theory Primer}\label{appendix-c-a-type-theory-primer}}

Like Appendix B, this appendix is scoped narrowly, to exactly what makes precise a handful of asides in the main text: the \texttt{xparse} argument specifiers from Chapter 6, the key-value interface design discussed in Chapter 7, and the document-class-as-type-constructor framing from Chapter 16. It is not an introduction to type theory as a foundation for mathematics or to the dependently typed proof assistants (Coq, Agda, Lean) that a reader might encounter elsewhere under this name. Nothing in the main chapters requires having read this appendix; a reader content to treat an \texttt{xparse} specifier as ``the pattern of arguments this macro accepts'' already has everything Chapters 6, 7, and 16 need.

\hypertarget{c.1-types-as-descriptions-of-shape}{%
\section{Types as Descriptions of Shape}\label{c.1-types-as-descriptions-of-shape}}

A type, in the narrow sense this appendix needs, is a description of the shape a piece of data is allowed to take: what it can be, before asking what its actual value is on any given occasion. Saying ``this argument has type integer'' says nothing about which integer; it says only that whatever value shows up there, it will be one of the things an integer can be, not a string, not a list, not something malformed. This is a useful description to attach to a macro's arguments for exactly the reason it is useful in an ordinary programming language: it lets an error be caught, or at least a mismatch be recognized, based on shape alone, before anyone has to reason about the specific value involved.

TeX itself has no type system in any of these senses; a macro argument is, at the primitive level, an unexamined sequence of tokens, and it is entirely the macro's own body that decides what to do with it, including whether to check its shape at all. Everything in this appendix is therefore describing a discipline layered on top of TeX by tools like \texttt{xparse} and \texttt{l3keys}, and by an author's own conventions, not a property TeX itself enforces the way a statically typed compiler would.

\hypertarget{c.2-simple-types-and-macro-signatures}{%
\section{Simple Types and Macro Signatures}\label{c.2-simple-types-and-macro-signatures}}

Chapter 6 introduced \texttt{xparse}'s argument specifiers, \texttt{m} for mandatory, \texttt{o} for optional, \texttt{s} for an optional star, \texttt{d} for a custom-delimited argument, combined into a specifier string describing a macro's full calling convention. Read through the lens of this appendix, a specifier string like \texttt{mo} is a simple type for the macro's argument list: it says the macro accepts exactly two arguments, the first always present, the second present or absent, and nothing about what values those arguments can meaningfully hold once supplied, only their presence, absence, and delimiting syntax.

This is deliberately a weak type system, in the sense that it checks shape (how many arguments, delimited how) without checking content (whether the first argument, once captured, is actually a sensible thing for the macro's body to do with). A macro declared with specifier \texttt{m} still receives whatever token sequence the caller wrote for that argument, with no guarantee it makes sense as, say, a length or a color name; any deeper checking has to be written into the macro's body explicitly, using the conditionals from Chapter 7 or an explicit format check. The value of the \texttt{xparse} specifier, even at this shallow level, is that it makes a macro's calling convention legible at the point of definition, \texttt{\textbackslash{}NewDocumentCommand\{\textbackslash{}keyword\}\{m\ o\}\{...\}} states the shape directly, rather than requiring a reader to infer the same information by reading through a \texttt{\textbackslash{}def}-based implementation that manually tests for and strips an optional bracket.

\hypertarget{c.3-records-and-key-value-interfaces}{%
\section{Records and Key-Value Interfaces}\label{c.3-records-and-key-value-interfaces}}

A key-value interface, covered practically in Chapter 7, corresponds to what a type theorist would call a record type: a structure with several independently named fields, each with its own type, where a value of the record type provides (or, for optional fields, may omit) a value for each field. \texttt{\textbackslash{}examplebox{[}boxed=false{]}\{...\}}, the worked example from Chapter 7, is supplying a partial record: a value for the \texttt{boxed} field, and, implicitly, defaults for \texttt{title} and \texttt{numbered}, the two fields left unspecified.

This framing clarifies something about good key-value interface design that Chapter 7 argued for on more informal grounds: every field of a well-designed record type has a sensible default, or is a genuinely required field the type does not permit omitting, and a record type's fields are, absent some explicit dependency between them, independently specifiable, which is exactly the composability property Chapter 7 asked a good key-value interface to have. A key-value interface that violates composability, where setting one key silently changes what another key means, is, in this language, not really a record type at all but something closer to a variant or tagged union pretending to be a record, and is harder to reason about, and to document, for exactly the reason a variant type disguised as a record would be harder to reason about in a conventional statically typed language.

\hypertarget{c.4-dependent-records-when-one-keys-validity-depends-on-another}{%
\section{Dependent Records: When One Key's Validity Depends on Another}\label{c.4-dependent-records-when-one-keys-validity-depends-on-another}}

Some key-value interfaces genuinely do need one key's valid values, or its very presence, to depend on another key's value: a \texttt{format} key that is only meaningful once a \texttt{type} key has been set to a specific value, or a \texttt{numbered} key whose valid range depends on which numbering scheme a separate \texttt{scheme} key selected. This is what a dependent record type describes: a record where a later field's type is itself computed from an earlier field's value, rather than being fixed in advance independent of everything else in the record.

\texttt{l3keys} supports this pattern directly, though not by that name, through key groups and conditional key definitions, where a key's own definition can inspect the current value of another key already processed earlier in the same call and behave differently depending on it. Recognizing this as ``the interface has a genuine dependency between two of its fields'' rather than merely ``this key's behavior is a little complicated'' is useful design guidance in its own right: a dependent key-value interface is harder to document, harder to validate, and harder for a caller to use correctly than an interface whose fields are all independent, so a dependency worth having is one earning its complexity, not one introduced merely because the implementation happened to make it easy to check one key's value while processing another.

\hypertarget{c.5-document-classes-as-type-constructors}{%
\section{Document Classes as Type Constructors}\label{c.5-document-classes-as-type-constructors}}

Chapter 16 described a custom document class as something that can require and structure a document's front-matter fields, title, author, affiliation, at the class level, raising an error at compile time if a required field is missing. In this appendix's terms, a document class defined this way is acting as a type constructor: \texttt{\textbackslash{}documentclass\{myresearchclass\}} does not merely load a set of macros, it instantiates a structured type, the class's own notion of ``a valid document of this class,'' with required fields that must be supplied (title, author) and optional fields with defaults (a subtitle, defaulting to none).

A class-level check that a required field is missing is, in this framing, a type error caught at compile time rather than a mistake a human proofreader has to catch by inspection, exactly the property Chapter 16 argued made this worth building into the class rather than leaving to convention. This is also why Chapter 16's advice to pin a shared class's version per book, so that an earlier book's build does not change when the shared class evolves, matters in the same way that changing a type's definition in a conventional codebase requires care: every existing ``value'' of that type, every book already built against an earlier version of the class, is implicitly relying on the type's earlier shape, and a change to what the type requires or permits can silently invalidate documents that were previously well-formed under the old definition.

\hypertarget{c.6-where-this-appendix-stops}{%
\section{Where This Appendix Stops}\label{c.6-where-this-appendix-stops}}

This appendix has covered exactly enough to read Chapter 6's \texttt{xparse} specifiers as lightweight signatures, Chapter 7's key-value interfaces as record types (dependent, where the interface genuinely needs it), and Chapter 16's document classes as type constructors enforcing a required shape at compile time. It has deliberately not covered the machinery a reader might associate with ``type theory'' from a more formal context: the Curry-Howard correspondence between types and propositions, dependent function types in the fuller sense used by proof assistants, or anything resembling a soundness proof for the informal type discipline described here, since TeX's macro system was never designed to support one and nothing in this book's practical guidance depends on one existing. A reader who wants that fuller treatment is better served by a dedicated type theory text, such as Pierce's \emph{Types and Programming Languages}~\cite{pierce2002types} for the programming-languages perspective this appendix has been borrowing vocabulary from, than by an appendix whose entire purpose was to stay exactly as large as three chapters' worth of asides required.
